\section{Methods}

We implement the Information--Cosserat framework using two complementary numerical approaches: a fast deterministic beam model for parameter sweeps and eigenanalysis, and a full three-dimensional Cosserat rod simulation for capturing large deformations and geometric nonlinearities.

\subsection{Deterministic IEC Beam Model}

To explore the mode selection mechanism (Eq.~\ref{eq:mode_selection}), we discretize the linearized beam equations using a finite difference scheme on a 1D domain $s \in [0, L]$. The rod is divided into $N=100$ segments. The information field $I(s)$ is mapped to local stiffness $E_i$ and rest curvature $\kappa_{i}$ at each node. We solve the resulting boundary value problem (BVP) using a standard shooting method (or sparse matrix solver for the eigenproblem). This allows rapid exploration of the $(\chi_\kappa, g)$ parameter space to identify regions where S-modes become the ground state.

\subsection{3D Cosserat Rod Implementation (PyElastica)}

For full 3D simulations, we utilize \texttt{PyElastica}~\cite{pyelastica_zenodo,gazzola2018forward}, an open-source Python implementation of Cosserat rod theory. Model parameters are listed in Table~\ref{tab:parameters}. The spine is modeled as a Cosserat rod with the following specifications:

\begin{itemize}
    \item \textbf{Discretization}: The rod is discretized into $n=50$--$100$ elements.
    \item \textbf{Information Field}: For spinal simulations, $I(s)$ is specified as a bimodal Gaussian distribution representing HOX-patterned identity:
    \begin{equation}
    I(s) = A_c \exp\left[-\frac{(s/L - s_c)^2}{2\sigma_c^2}\right] + A_l \exp\left[-\frac{(s/L - s_l)^2}{2\sigma_l^2}\right] + I_0,
    \label{eq:info_field_spinal}
    \end{equation}
    where $A_c = 0.5$, $s_c = 0.80$, $\sigma_c = 0.08$ (cervical lordosis), $A_l = 0.7$, $s_l = 0.25$, $\sigma_l = 0.10$ (lumbar lordosis), and $I_0 = 0.3$ (baseline). For scoliosis perturbations, a lateral asymmetry is added: $I(s) \to I(s) + \varepsilon_{\mathrm{asym}} \exp[-(s/L - 0.6)^2/(2 \cdot 0.08^2)]$ with $\varepsilon_{\mathrm{asym}} = 0.01$--$0.05$.
    \item \textbf{IEC Coupling}: We implemented a custom callback in \texttt{PyElastica} that updates the local rest curvature vector $\bm{\kappa}^0(s)$ and bending stiffness matrix $\mathbf{B}(s)$ at each time step (or initialization) based on the information field $I(s)$.
    \item \textbf{Boundary Conditions}: The rod is clamped at the base (sacrum) and free at the top (cranium), simulating a cantilever column under gravity. For specific validation cases, clamped-clamped conditions are used.
    \item \textbf{Gravitational Loading}: Gravity is applied as a uniform body force $\mathbf{f} = \rho A \mathbf{g}$.
    \item \textbf{Damping}: To find static equilibrium configurations, we apply external damping ($\nu \sim 0.1$--$1.0$ s$^{-1}$) and integrate the dynamic equations until the kinetic energy dissipates ($v_{\max} < 10^{-6}$ m/s).
\end{itemize}

The source code for the IEC-modified Cosserat solver is available in the \texttt{spinalmodes} Python package (see Data Availability).

\subsection{Parameter Sweeps and Mode Classification}

We perform systematic parameter sweeps over the coupling strength $\chi_\kappa$ (range $[0, 0.1]$) and gravitational acceleration $g$ (range $[0.01, 1.0]$ $g_{\mathrm{Earth}}$). For each simulation, we compute the equilibrium shape and evaluate the following metrics:
1.  \textbf{Geodesic Deviation} $\widehat{D}_{\mathrm{geo}}$: Quantifies the difference between the realized shape and the gravity-only geodesic.
2.  \textbf{Lateral Deviation} $S_{\mathrm{lat}}$: Measures symmetry breaking in the coronal plane.
3.  \textbf{Cobb Angle}: Standard clinical measure for scoliotic curves.

Regimes are classified based on the normalized geodesic deviation $\widehat{D}_{\mathrm{geo}}$. We define the \emph{gravity-dominated} regime ($\widehat{D}_{\mathrm{geo}} < 0.1$) as the region where gravitational potential energy dominates, leading to monotonic sag. The \emph{cooperative} regime ($0.1 \leq \widehat{D}_{\mathrm{geo}} \leq 0.3$) corresponds to the range where information and gravity balance to produce a stable, counter-curved S-shape. The \emph{information-dominated} regime ($\widehat{D}_{\mathrm{geo}} > 0.3$) is reached when information-driven metric distortions become the primary determinant of geometry, which, as we show, also coincides with the onset of symmetry-breaking instabilities (scoliosis-like modes). These thresholds were determined by analyzing the bifurcation of the lowest eigenvalue $\lambda_0$ and the emergence of non-monotonic curvature profiles.

\subsection{Multi-Scale Protein Mechanics Pipeline}

To ground the macroscopic parameters in molecular reality, we developed a multi-scale pipeline integrating AlphaFold predictions with structural energetics.

\paragraph{Protein Structure Prediction.} Structures were retrieved from the AlphaFold Protein Structure Database (v4) where available and predicted de novo otherwise. For each protein, we extracted pLDDT, domain boundaries, radius of gyration, and anisotropy ratio. Proteins with extensive disordered regions (pLDDT $< 70$) were retained but flagged for larger mechanical uncertainty.

\paragraph{Steered Molecular Dynamics (SMD).} SMD was run in GROMACS 2023.3 with CHARMM36m and TIP3P water. Systems were energy-minimized, equilibrated under NVT then NPT ensembles (310 K, 1 bar), and pulled with a harmonic spring along the principal molecular axis ($v_{\mathrm{pull}}=0.01$ nm/ps, $k=1000$ kJ/(mol$\cdot$nm$^2$)). Force-extension curves were fitted in the low-strain regime to estimate stiffness. Type I collagen served as a calibration anchor ($k \approx 420 \pm 65$ pN/nm).

\paragraph{Fibril Assembly and Coarse-Grained Mechanics.} Atomistic outputs were coarse-grained with Martini 3 and assembled into quasi-hexagonal fibrils with 67 nm D-band staggering. Proteoglycan contributions were added through an osmotic term based on Donnan equilibrium. Composite fibril-matrix behavior was estimated by mixture rules with composition-dependent volume fractions.

\paragraph{Tissue Homogenization.} Tissue-level elasticity was computed as:
\begin{equation}
\mathbf{C}_{\mathrm{tissue}}(s) = \sum_i \phi_i(s)\,\mathbf{R}(\theta_i)\,\mathbf{C}_{\mathrm{fibril},i}\,\mathbf{R}(\theta_i)^T
\end{equation}
where $\phi_i(s)$ are volume fractions from histology, $\theta_i$ are orientation angles from polarized imaging, and $\mathbf{C}_{\mathrm{fibril},i}$ are constituent stiffness tensors. This parameterization yielded regional variations: $E_{\mathrm{cervical}}=1.15 \pm 0.18$ GPa, $E_{\mathrm{thoracic}}=0.32 \pm 0.07$ GPa, and $E_{\mathrm{lumbar}}=0.78 \pm 0.12$ GPa, with predicted-versus-measured agreement $R^2=0.81$ across 12 vertebral levels.

\subsection{Protein Dynamics and Energy Modeling}

We simulated the temporal competition between mechanosensory and growth protein pools using coupled kinetics:
\begin{align}
\frac{dP_{\mathrm{mech}}}{dt} &= \alpha_m(E_{\mathrm{mech}}) - \delta_m P_{\mathrm{mech}} \\
\frac{dP_{\mathrm{grow}}}{dt} &= \alpha_g(E_{\mathrm{grow}}) - \delta_g P_{\mathrm{grow}}
\end{align}
subject to $\alpha_m + \alpha_g \leq E_{\mathrm{total}} - E_{\mathrm{metabolic}}$. Degradation constants were set to $\delta_m=0.058$ h$^{-1}$ and $\delta_g=0.001$ h$^{-1}$. Growth forcing followed fitted adolescent velocity curves. At each time step, ATP-constrained synthesis rates generated the instability ratio $\Lambda(t)$.

\subsection{Tensor Elasticity and Vector Field Analysis}

The alignment map $\alpha(s)$ was computed nodewise from information-gradient and stress vectors. Vector mismatch regions were defined as contiguous segments with $\alpha<0.5$, and severe mismatch as $\alpha<0.3$. These thresholds were applied consistently in model and imaging analyses.

\subsection{Numerical convergence and validation}

A convergence study was performed by varying the number of elements $n$ from $10$ to $200$. We found that the normalized geodesic deviation $\widehat{D}_{\mathrm{geo}}$ converges to within 1\% for $n \geq 50$ (see Supplementary Fig.~S1). The numerical implementation was validated against analytical solutions for small-deflection Euler-Bernoulli beams. The \texttt{PyElastica} implementation was further verified by reproducing standard buckling and hanging chain benchmarks, with error metrics $||\mathbf{r}_{num} - \mathbf{r}_{ana}||/L < 10^{-4}$.
